\section{Volumes e Bind Mounts}

\begin{frame}
  \frametitle{O Problema da Persistência}
  Quando um container é removido, \textbf{todas as alterações de arquivos são perdidas}.

  \vspace{0.5em}
  Exemplo: reiniciar um container de banco de dados = banco vazio.

  \vspace{0.8em}
  O Docker oferece duas soluções:
  \begin{itemize}
    \item \textbf{Volumes} — gerenciados pelo Docker
    \item \textbf{Bind Mounts} — mapeiam diretório do host
  \end{itemize}
\end{frame}

\begin{frame}
  \frametitle{Volumes}
  Mecanismo de armazenamento \textbf{gerenciado pelo Docker} que persiste dados independentemente do container.

  \vspace{0.5em}
  \texttt{docker volume create meus-dados}

  \vspace{0.3em}
  \texttt{docker run -d -v meus-dados:/app/data minha-imagem}

  \vspace{0.8em}
  \begin{itemize}
    \item Se o volume não existir, Docker cria automaticamente
    \item Dados sobrevivem à remoção do container
    \item Podem ser compartilhados entre múltiplos containers
  \end{itemize}
\end{frame}

\begin{frame}
  \frametitle{Gerenciando Volumes}
  \texttt{docker volume ls} \hspace{2.5em} — listar volumes

  \vspace{0.3em}
  \texttt{docker volume rm <nome>} \hspace{0.3em} — remover um volume

  \vspace{0.3em}
  \texttt{docker volume prune} \hspace{1.3em} — remover não utilizados

  \vspace{0.5em}
  \begin{alertblock}{Atenção}
    Um volume só pode ser removido se não estiver conectado a nenhum container.
  \end{alertblock}
\end{frame}

\begin{frame}[fragile]
  \frametitle{Exemplo: PostgreSQL com Volume}
  \small
\begin{verbatim}
# Iniciar com volume
docker run --name=db -e POSTGRES_PASSWORD=secret \
  -d -v postgres_data:/var/lib/postgresql postgres:18

# Inserir dados
docker exec -ti db psql -U postgres
CREATE TABLE tasks (id SERIAL, description VARCHAR(100));
INSERT INTO tasks (description) VALUES ('Finish work');
\q

# Remover e recriar
docker stop db && docker rm db
docker run --name=new-db -d \
  -v postgres_data:/var/lib/postgresql postgres:18

# Dados ainda estão lá!
docker exec -ti new-db psql -U postgres \
  -c "SELECT * FROM tasks"
\end{verbatim}
\end{frame}

\begin{frame}
  \frametitle{Bind Mounts}
  Mapeiam \textbf{diretamente} um diretório do host para dentro do container.

  \vspace{0.5em}
  \texttt{docker run -v ./local:/app/config nginx}

  \vspace{0.3em}
  \texttt{docker run --mount type=bind,source=./,target=/app nginx}

  \vspace{0.8em}
  \begin{itemize}
    \item Ideais para \textbf{desenvolvimento} (acesso em tempo real)
    \item Com \texttt{-v}: diretório criado automaticamente se não existir
    \item Com \texttt{--mount}: erro se não existir (recomendado pelo Docker)
  \end{itemize}

  \vspace{0.5em}
  Permissões: \texttt{:ro} (somente leitura) ou \texttt{:rw} (leitura/escrita, padrão)
\end{frame}

\begin{frame}
  \frametitle{Volume vs Bind Mount}
  \begin{center}
  \begin{tabular}{|l|l|}
  \hline
  \textbf{Volume} & \textbf{Bind Mount} \\
  \hline
  Gerenciado pelo Docker & Caminho direto no host \\
  Dados de produção & Desenvolvimento \\
  Compartilhável & Sincronização host-container \\
  Backup facilitado & Acesso direto aos arquivos \\
  \hline
  \end{tabular}
  \end{center}
\end{frame}
